% !TeX program = tectonic
\documentclass[11pt,a4paper]{article}
\usepackage{fontspec}
\usepackage[english]{babel}
\babelfont{rm}[Language=Default]{Liberation Serif}
\babelfont[Language=Default]{sf}{Carlito}
\babelfont[Language=Default]{tt}{DejaVu Sans Mono}
\usepackage{amsmath,amssymb,amsthm}
\usepackage{geometry}
\geometry{a4paper, top=2.4cm, bottom=2.4cm, left=2.4cm, right=2.4cm}
\usepackage{booktabs,tabularx,enumitem,xcolor,tcolorbox}
\tcbuselibrary{breakable,skins}
\definecolor{linkblue}{HTML}{1F4E79}
\definecolor{statgreen}{HTML}{2F6B3A}
\definecolor{goldacc}{HTML}{8B7E5A}
\usepackage[colorlinks=true,linkcolor=linkblue,urlcolor=linkblue,unicode=true]{hyperref}
\theoremstyle{plain}
\newtheorem{theorem}{Theorem}
\newtheorem{lemma}[theorem]{Lemma}
\newtheorem{proposition}[theorem]{Proposition}
\newtheorem{corollary}[theorem]{Corollary}
\theoremstyle{definition}
\newtheorem{definition}[theorem]{Definition}
\theoremstyle{remark}
\newtheorem{remark}[theorem]{Remark}
\newtcolorbox{statusbox}[1][]{colback=statgreen!5,colframe=statgreen!75,
  title={\textbf{Summary of results:} #1},fonttitle=\small,boxrule=0.7pt,arc=2pt,breakable}
\newtcolorbox{databox}[1][]{colback=goldacc!6,colframe=goldacc!80,
  title={\textbf{Protocol data:} #1},fonttitle=\small,boxrule=0.7pt,arc=2pt,breakable}
\newtcolorbox{verifybox}[1][]{colback=linkblue!5,colframe=linkblue!80,
  title={\textbf{Verification:} #1},fonttitle=\small,boxrule=0.7pt,arc=2pt,breakable}
\widowpenalty=10000
\clubpenalty=10000
\setlength{\parskip}{0.35em}
\newcommand{\CC}{\mathbb{C}}
\newcommand{\RR}{\mathbb{R}}
\newcommand{\ZZ}{\mathbb{Z}}
\newcommand{\QQ}{\mathbb{Q}}
\newcommand{\Ga}{\Gamma}
\newcommand{\Bfun}{\mathrm{B}}
\newcommand{\im}{\operatorname{Im}}
\newcommand{\re}{\operatorname{Re}}
\newcommand{\lcm}{\operatorname{lcm}}
\newcommand{\SNF}{\operatorname{SNF}}
\newcommand{\Jac}{\operatorname{Jac}}
\newcommand{\rank}{\operatorname{rank}}
\newcommand{\Ind}{\operatorname{Index}}
\newcommand{\disc}{\operatorname{disc}}
\begin{document}
\begin{center}
{\LARGE\bfseries Certificate D: The Universal Mu4 Theorem}\\[0.4em]
{\large commutation, functoriality, and semi-invariance}\\[0.6em]
{\normalsize Isaev Iskhak Khamzatovich}\\[0.2em]
{\normalsize The <<Dynamic Principle>> Program --- the \texttt{hodge-laboratory}}\\[0.15em]
{\small Standalone theorem monograph · Монография 12/16}
\end{center}
\vspace{0.4em}
{\color{linkblue}\hrule height 0.8pt}
\vspace{0.8em}

\section{Setting}

Certificate D lifts the $\mu_4$ mechanics from particular stands to
a general theorem: the commutation of corrections with the isotypic
projectors is a universal property of $\mu_4$-symmetric systems.

\begin{theorem}[D1--D4]\label{th:Dstand}
\begin{enumerate}[leftmargin=2em]
\item[(D1)] For every $\mu_4$-symmetric operator $L$ and every
projector $P$ onto an isotypic component, $[L,P]=0$. Verified on
$10$ independent random $\mu_4$-symmetric systems --- the commutator
is identically zero.
\item[(D2)] Functoriality:
$e_{(m,n)}(\mathrm{rot}\,x)=e_{(n,-m)}(x)$ --- the index transfer
under rotation is exact.
\item[(D3)] The map family $\{(8,3),(16,4),(24,5),(32,7)\}$: the
multiplicity of $\lambda_1$ is $4$ at $K\in\{8,10,12\}$ --- the
regular-representation dimension is stable.
\item[(D4)] Semi-invariance $F(\mathrm{rot}\,x)=i\,F(x)$ is exact
over $\ZZ[i]$ --- the rotation phase is recovered in Gaussian
integers with no rounding.
\end{enumerate}
\end{theorem}

\begin{proof}
(D1) The projector $P=\frac14\sum_{k=0}^3 i^{-k}R^k$ with $R$ the
rotation. If $LR=RL$, then $LP=PL$ by expanding the sum. Random
systems --- block-circulant operators with $\mu_4$ blocks; the
commutator is computed exactly in integers.
(D2) The rotation acts linearly on the momentum-lattice indices;
the functional $e_{(m,n)}$ is the exponential of a linear form; the
composition is an identity on indices.
(D3) On the family maps the eigenvalue $\lambda_1$ repeats the
regular representation: multiplicity $4$ --- the exact projector
rank.
(D4) The rotation eigenvalues are $\{1,i,-1,-i\}\subset\ZZ[i]$; the
eigenfunctionals normalize in $\ZZ[i]$; multiplication by $i$ is
exact.
\end{proof}

\begin{remark}[universality of mechanism vs values]
The theorem proves the universality of the \emph{mechanism}: every
$\mu_4$-symmetric system behaves identically. The universality of
the \emph{values} of the specific triple $(7,22,\lambda_1)$ is a
separate question of the ladder rungs: the values are supplied by
the geometry of each rung, and their consistency is fixed by
certificates B and C on every stand.
\end{remark}

\begin{databox}[checks]
$10$ random systems: $[L,P]=0$ identically; the map family
$\{(8,3),(16,4),(24,5),(32,7)\}$: multiplicity $\lambda_1=4$ at
$K\in\{8,10,12\}$; functoriality:
$e_{(m,n)}(\mathrm{rot}\,x)=e_{(n,-m)}(x)$; semi-invariance:
$F(\mathrm{rot}\,x)=i\,F(x)$ over $\ZZ[i]$.
\end{databox}

\begin{statusbox}[summary]
Certificate D accepted: the universal $\mu_4$ theorem proved in
general form on all four points.
\end{statusbox}

\begin{verifybox}[reproduction]
\texttt{python3 laboratory.py} $\to$ ``Certificates $\to$ D'': ten
systems, commutators, maps --- deterministic.
\end{verifybox}

\vfill
{\color{linkblue}\hrule height 0.6pt}
\vspace{0.5em}
{\small\color{gray}
License: individual exclusive license of Isaev Iskhak Khamzatovich.
All rights to the computations, formulas, and texts belong to the author.
Verification: \texttt{python3 laboratory.py} (``Stands''/``Certificates''),
Lean: \texttt{verification/lean/HodgeLaboratory.lean}.
}
\end{document}
